\documentclass[../../../include/open-logic-section]{subfiles}

\begin{document}

\olfileid{sth}{ordinals}{basic}
\olsection{ક્રમવાચક સંખ્યાઓના મૂળભૂત ગુણધર્મો}

આપણે જોયું કે શરૂઆતની કેટલીક ક્રમવાચક સંખ્યાઓ
પ્રાકૃતિક સંખ્યાઓ જ છે. ક્રમવાચક સંખ્યાઓનો સિદ્ધાંત
વિકસાવવાનું મુખ્ય કારણ પ્રાકૃતિક સંખ્યાઓ પર સાચા
અનુમાનના સિદ્ધાંતને આગળ વિસ્તારવાનું છે. પ્રાથમિક
પરિણામોની શ્રેણીથી આપણે ત્યાં સુધી પહોંચીશું.

\begin{lem}\ollabel{ordmemberord}
ક્રમવાચક સંખ્યાનો દરેક !!{element} પણ ક્રમવાચક
સંખ્યા છે.
\end{lem}

\begin{proof}
$\alpha$ ક્રમવાચક સંખ્યા હોય અને $b \in \alpha$ હોય.
$\alpha$ પરંપરિત હોવાથી $b \subseteq \alpha$.
તેથી $\in$ એ $b$ને સુક્રમિત કરે છે, કારણ કે
$\in$ એ $\alpha$ને સુક્રમિત કરે છે.

$b$ પરંપરિત છે તે જોવા માટે ધારો કે $x \in c \in b$.
$b
\subseteq \alpha$ હોવાથી $c \in \alpha$.
ફરી $\alpha$ પરંપરિત હોવાથી $c \subseteq
\alpha$, એટલે $x \in \alpha$.
આમ $x, c, b \in \alpha$.
પણ $\in$ એ $\alpha$ને સુક્રમિત કરે છે, તેથી
\olref[sth][ordinals][wo]{wo:strictorder} મુજબ
$\alpha$ પરનો $\in$ સંબંધ પરંપરિત છે.
એટલે $x
\in c \in b$ હોવાથી $x \in b$.
સામાન્ય રીતે, $c \subseteq b$.
\end{proof}

\begin{cor}\ollabel{ordissetofsmallerord}
કોઈ પણ ક્રમવાચક સંખ્યા~$\alpha$ માટે
$\alpha = \Setabs{\beta \in \alpha}{\beta \text{ ક્રમવાચક સંખ્યા છે}}$.
\end{cor}

\begin{proof}
\olref{ordmemberord}થી સીધું મળે છે.
\end{proof}

આગળનાં બે મુખ્ય પરિણામો,
\olref{ordinductionschema} અને
\olref{ordtrichotomy}, આશરે એ કહે છે કે
ક્રમવાચક સંખ્યાઓ પોતે સભ્યતા-સંબંધથી સુક્રમિત છે:

\begin{thm}[અનંતાતીત અનુમાન]\ollabel{ordinductionschema}
કોઈ પણ સૂત્ર $\phi(x)$ માટે:
\[
	\text{જો }\exists \alpha \phi(\alpha)\text{, તો }\exists \alpha(\phi(\alpha)
	\land  (\forall \beta \in \alpha) \lnot \phi(\beta))
\]
અહીં દર્શાવેલા પરિમાણકો પરોક્ષ રીતે ક્રમવાચક
સંખ્યાઓ સુધી મર્યાદિત છે.
\end{thm}

\begin{proof}
%\olref{ordinductionschema1}. 
ધારો કે કોઈ ક્રમવાચક સંખ્યા $\alpha$ માટે
$\phi(\alpha)$ છે. જો
$ (\forall \beta
\in \alpha) \lnot \phi(\beta)$ હોય, તો કામ પૂરું.
નહિતર $\alpha$ના એવા સભ્યોમાંથી, જેમના માટે
$\phi$ સાચું છે, $\in$ મુજબ લઘુતમ !!{element}
પસંદ કરો. આપેલો ગણ ક્રમવાચક સંખ્યા હોવાથી
આવો સભ્ય છે અને તે \olref{ordmemberord}
મુજબ ક્રમવાચક સંખ્યા પણ છે.
\sourcecorrection{OLMD-116}{મૂળમાં માત્ર $\alpha$ને
કોઈ $\in$-લઘુતમ સભ્ય છે “જે $\phi$ છે”
એમ કહેવાયું છે. જરૂરી પસંદગી આખા $\alpha$માંથી
નહીં, પણ $\phi$ સંતોષતા તેના સભ્યોના અરિક્ત
ઉપગણમાંથી કરવી પડે: પહેલેથી $\phi(\alpha)$
હોય અને તે ઉપગણ ખાલી હોય તો $\alpha$ સાક્ષી છે;
નહિતર ઉપગણનો લઘુતમ સભ્ય સાક્ષી છે.}
%	\olref{ordinductionschema2}. Suppose there is some ordinal
%	$\gamma$ such that $\lnot\phi(\gamma)$. Then by
%	\olref{ordinductionschema1}, there is an $\in$-minimal ordinal
%	$\alpha$ for which $\lnot\phi(\alpha)$. So $(\forall \beta <
%	\alpha) \phi(\beta)$, rendering the antecedent of the conditional
%	false.
\end{proof}
\noindent
નોંધો કે \olref{ordinductionschema}ને સમકક્ષ રીતે
આ પ્રરૂપ તરીકે પણ વ્યક્ત કરી શકાય:
\[
\text{જો }\forall \alpha((\forall \beta \in \alpha)\phi(\beta) \lif
\phi(\alpha))\text{, તો }\forall \alpha\phi(\alpha)
\]
\olref{ordinductionschema}માં
$\lnot\phi(\alpha)$ મૂકીને અને પછી પ્રાથમિક
તાર્કિક ફેરફારો કરીને આ મળે છે.

\begin{thm}[ત્રિવિકલ્પતા]\ollabel{ordtrichotomy}
કોઈ પણ ક્રમવાચક સંખ્યાઓ $\alpha$ અને $\beta$ માટે
$\alpha \in \beta \lor \alpha = \beta \lor \beta \in \alpha$.
\end{thm}

\begin{proof}
સાબિતી બે વાર અનુમાનથી છે, એટલે કે
\olref{ordinductionschema} બે વાર વાપરીશું.
$x$ એ $y$ સાથે \emph{તુલનીય} છે એમ કહીએ
જો અને ફક્ત જો
$x \in y \lor x = y \lor y \in x$.

અનુમાન માટે ધારો કે~$\alpha$માંની દરેક
ક્રમવાચક સંખ્યા \emph{દરેક} ક્રમવાચક
સંખ્યા સાથે તુલનીય છે.
બીજા અનુમાન માટે ધારો કે
$\alpha$ એ~$\beta$માંની દરેક ક્રમવાચક
સંખ્યા સાથે તુલનીય છે.
આપણે $\alpha$ અને~$\beta$ તુલનીય છે તે
બતાવીશું. પછી~$\beta$ પરના અનુમાનથી
$\alpha$ દરેક ક્રમવાચક સંખ્યા સાથે
તુલનીય છે; અને ત્યાર બાદ~$\alpha$ પરના
અનુમાનથી \emph{દરેક} ક્રમવાચક સંખ્યા
\emph{દરેક} ક્રમવાચક સંખ્યા સાથે તુલનીય
છે, જેમ જોઈએ છે.
$\alpha \notin \beta$ અને
$\beta \notin
\alpha$ ધારીને $\alpha = \beta$ બતાવવું પૂરતું છે.

$\alpha \subseteq \beta$ બતાવવા કોઈ
$\gamma \in \alpha$ લો; \olref{ordmemberord}
મુજબ તે ક્રમવાચક સંખ્યા છે.
તેથી પહેલી અનુમાન-ધારણા મુજબ
$\gamma$ એ $\beta$ સાથે તુલનીય છે.
પણ જો $\gamma
= \beta$ અથવા $\beta \in \gamma$ હોય,
તો (જરૂર પડે ત્યાં $\alpha$ પરંપરિત છે તે
વાપરીને) $\beta \in \alpha$ થશે, જે
આપણી ધારણાની વિરુદ્ધ છે; તેથી
$\gamma \in \beta$. સામાન્ય રીતે,
$\alpha \subseteq
\beta$.

એ જ પ્રકારની દલીલ બીજી અનુમાન-ધારણા
વાપરીને $\beta \subseteq \alpha$ બતાવે છે.
તેથી $\alpha = \beta$.
\end{proof}\noindent તેથી આપણે ક્યારેક
$\alpha \in \beta$ને બદલે $\alpha <\beta$
લખીશું, કારણ કે $\in$ અહીં ક્રમસંબંધ
જેવો વ્યવહાર કરે છે. પરિચિતપણું અને
$\alpha \in
\beta \lor \alpha = \beta$ કરતાં
$\alpha \leq \beta$ લખવું સરળ હોવું
એ સિવાય આ માટે કોઈ ગહન કારણ નથી.\footnote{
આપણે $\alpha
\mathrel{\underline{\in}} \beta$ લખી શકીએ,
પણ તે તદ્દન અપ્રચલિત સંકેત થશે.}

આ છેલ્લાં પરિણામોનાં બે ટૂંકાં ફલિતાર્થો
આ છે; પહેલું આપણા નવા સંકેતનો ઉપયોગ કરે છે:

\begin{cor}\ollabel{ordordered}
જો $\exists \alpha\phi(\alpha)$ હોય, તો
$\exists \alpha(\phi(\alpha)
\land \forall \beta(\phi(\beta) \lif \alpha \leq \beta))$.
વધુમાં, કોઈ પણ ક્રમવાચક સંખ્યાઓ
$\alpha, \beta, \gamma$ માટે
$\alpha
\notin \alpha$ અને
$\alpha \in \beta \in \gamma \lif \alpha \in
\gamma$ બંને સાચાં છે.
\end{cor}

\begin{proof}
\olref[wo]{wo:strictorder} જેવી જ દલીલ.
\end{proof}

\begin{prob}
\olref[sth][ordinals][basic]{ordtrichotomy}ની
સાબિતીમાં “એ જ પ્રકારની દલીલ” પૂરી કરો.
\end{prob}

\begin{cor}\ollabel{corordtransitiveord}
$A$ ક્રમવાચક સંખ્યા હોય જો અને ફક્ત જો
$A$ એ ક્રમવાચક સંખ્યાઓનો પરંપરિત ગણ હોય.
\end{cor}

\begin{proof}
\emph{ડાબેથી જમણે.} \olref{ordmemberord} મુજબ.
\emph{જમણેથી ડાબે.}
જો $A$ ક્રમવાચક સંખ્યાઓનો પરંપરિત
ગણ હોય, તો
\olref{ordinductionschema} અને
\olref{ordtrichotomy} મુજબ
$\in$ એ $A$ને સુક્રમિત કરે છે.
\end{proof}

હમણાં આપણે \olref{ordinductionschema}
અને \olref{ordtrichotomy}નો સાર એ રીતે
કહ્યો કે $\in$ ક્રમવાચક સંખ્યાઓને
સુક્રમિત કરે છે. જોકે આગળના પરિણામને
લીધે આવા દાવા અંગે આપણે \emph{ખૂબ સાવચેત}
રહેવું જોઈએ:

\begin{thm}[બુરાલી–ફૉર્ટીનો વિરોધાભાસ]\ollabel{buraliforti}
બધી ક્રમવાચક સંખ્યાઓનો ગણ નથી.
\end{thm}

\begin{proof}
વિરોધાભાસ માટે ધારો કે $O$ બધી
ક્રમવાચક સંખ્યાઓનો ગણ છે.
જો $\alpha \in
\beta \in O$ હોય, તો
\olref{ordmemberord} મુજબ $\alpha$
ક્રમવાચક સંખ્યા છે, તેથી
$\alpha \in O$.
આમ $O$ પરંપરિત છે અને તેથી
\olref{corordtransitiveord} મુજબ
$O$ ક્રમવાચક સંખ્યા છે.
એટલે $O \in O$, જે
\olref{ordordered} સાથે વિરોધાભાસ કરે છે.
\end{proof}
\noindent
આ પરિણામનું નામ \citeauthor{Burali-Forti1897}
પરથી પડ્યું છે. પરંતુ 1899માં ડેડેકિન્ડને
લખેલા પત્રમાં કૅન્ટૉરે બધી ક્રમવાચક
સંખ્યાઓનો ગણ માનતાં ઊભો થતો
\emph{વિરોધાભાસ} પહેલી વાર સ્પષ્ટ રીતે
જોયો હતો. વાન હેયેનૉર્ટ સમજાવે છે તેમ:
\begin{quote}
  બુરાલી–ફૉર્ટીએ પોતે વિરોધાભાસને
  \emph{વિરોધાભાસ વડે સાબિતી}થી
  એવું પરિણામ મળ્યું હોવાનું માન્યું કે
  ક્રમવાચક સંખ્યાઓનો સ્વાભાવિક ક્રમ
  માત્ર આંશિક ક્રમ છે.
  \citep[p.~105]{Heijenoort1967}
\end{quote}
બુરાલી–ફૉર્ટીની ભૂલને બાજુએ રાખીને
અત્યાર સુધીની વાત આ રીતે કહી શકીએ:
ક્રમવાચક સંખ્યાઓ એવા ગણો છે જે
અલગ અલગ રીતે સભ્યતા-સંબંધથી
સુક્રમિત છે, અને સામૂહિક રીતે પણ
સભ્યતા-સંબંધથી સુક્રમિત છે
(પણ એ બધી મળીને ગણ બનતી નથી).

છેલ્લે, ક્રમવાચક સંખ્યાઓના
\olref{ordinductionschema} અને
\olref{ordtrichotomy}માંથી મળતાં
કેટલાક વધુ મૂળભૂત ગુણધર્મો આ છે.

\begin{prop}
ક્રમવાચક સંખ્યાઓનો કોઈ પણ કડક
ઊતરતો અનુક્રમ સાન્ત હોય છે.
\end{prop}

\begin{proof}
ક્રમવાચક સંખ્યાઓના કોઈ અનંત કડક
ઊતરતા અનુક્રમ
$\alpha_0 > \alpha_1 > \alpha_2 > \ldots$માં
$<$ મુજબ ન્યૂનતમ સભ્ય નથી;
આ \olref{ordinductionschema}ની
વિરુદ્ધ છે.
\end{proof}

\begin{prop}\ollabel{ordinalsaresubsets}
કોઈ પણ ક્રમવાચક સંખ્યાઓ
$\alpha, \beta$ માટે
$\alpha \subseteq \beta \lor \beta \subseteq \alpha$.
\end{prop}

\begin{proof}
જો $\alpha \in \beta$, તો $\beta$
પરંપરિત હોવાથી $\alpha \subseteq \beta$.
એ જ રીતે, જો $\beta \in \alpha$,
તો $\beta \subseteq
\alpha$. અને જો $\alpha = \beta$,
તો $\alpha \subseteq \beta$ અને
$\beta \subseteq \alpha$.
તેથી \olref{ordtrichotomy} મુજબ
દલીલ પૂરી થાય છે.
\end{proof}

\begin{prop}\ollabel{ordisoidentity}
કોઈ પણ ક્રમવાચક સંખ્યાઓ
$\alpha, \beta$ માટે
$\alpha = \beta$ તો અને ફક્ત તો
$\ordeq{\alpha}{\beta}$.
\end{prop}

\begin{proof}
ક્રમવાચક સંખ્યાઓ સુક્રમો છે;
તેથી ત્રિવિકલ્પતા
(\olref[basic]{ordtrichotomy}) અને
\olref[iso]{wellordnotinitial}થી
આ સીધું મળે છે.
\end{proof}

\begin{prob}\ollabel{probunionordinalsordinal}
જો $X$નો દરેક સભ્ય ક્રમવાચક
સંખ્યા હોય, તો $\bigcup X$
ક્રમવાચક સંખ્યા છે તે સાબિત કરો.
\end{prob}

\end{document}
